\documentclass[12pt,a4paper]{article}

\usepackage{amsmath,amssymb,amsthm}
\usepackage{mathtools}
\usepackage{tikz}
\usetikzlibrary{arrows.meta,positioning,cd}
\usepackage{enumitem}
\usepackage[margin=2.5cm]{geometry}
\usepackage{hyperref}
\usepackage{booktabs}

% Theorem environments
\newtheorem{theorem}{Theorem}[section]
\newtheorem{proposition}[theorem]{Proposition}
\newtheorem{lemma}[theorem]{Lemma}
\newtheorem{corollary}[theorem]{Corollary}
\theoremstyle{definition}
\newtheorem{definition}[theorem]{Definition}
\newtheorem{example}[theorem]{Example}
\theoremstyle{remark}
\newtheorem{remark}[theorem]{Remark}

% Shortcuts
\newcommand{\C}{\mathbb{C}}
\newcommand{\Z}{\mathbb{Z}}
\newcommand{\N}{\mathbb{N}}
\newcommand{\F}{\mathbb{F}}
\newcommand{\gl}{\mathfrak{gl}}
\newcommand{\fU}{\mathfrak{U}}
\newcommand{\fD}{\mathfrak{D}}
\newcommand{\fL}{\mathfrak{L}}
\newcommand{\End}{\operatorname{End}}
\newcommand{\Hom}{\operatorname{Hom}}
\newcommand{\Ind}{\operatorname{Ind}}
\newcommand{\Res}{\operatorname{Res}}
\newcommand{\Stab}{\operatorname{Stab}}
\newcommand{\GL}{\mathrm{GL}}
\newcommand{\Rep}{\mathcal{R}\mathrm{ep}}
\newcommand{\Vect}{\mathcal{V}\!\mathit{ect}}
\newcommand{\Kl}{\mathrm{Kl}}
\newcommand{\EM}{\mathrm{EM}}
\newcommand{\Sab}{\mathrm{Sab}}
\newcommand{\Hecke}{\mathcal{H}}

\title{The Mackey Formula as Monad Decomposition}
\author{Notes}
\date{}

\begin{document}
\maketitle

\begin{abstract}
We give a detailed account of the Mackey formula through two worked
examples: $S_n$ inside $S_{n+1}$ (two double cosets, the Heisenberg
relation) and $\GL_n(\F_q)$ inside $\GL_{n+1}(\F_q)$ ($q^2+2q-2$
double cosets, new field-arithmetic channels). We explain the
monad $\Ind \circ \Res$, whose Mackey decomposition recovers the
Hecke algebra as Kleisli endomorphisms, and connect to Sabidussi
coset graphs and Fomin's local rules.
\end{abstract}

\tableofcontents

%% ====================================================================
\section{The Mackey formula}
\label{sec:mackey}
%% ====================================================================

\subsection{Double cosets}

Let $G$ be a finite group and $H \leqslant G$ a subgroup. The set of
\emph{double cosets} $H \backslash G / H$ partitions $G$ into disjoint
subsets $HgH = \{hgk : h, k \in H\}$.

The key geometric idea: every conjugate $gHg^{-1}$ is another copy of
$H$ inside $G$. The double coset $HgH$ measures \emph{how this
conjugate copy sits relative to the original}, up to the symmetry of
$H$ itself. Two elements $g, g'$ give the same double coset if and
only if the ``relative position'' of $gHg^{-1}$ and $H$, as seen by
$H$, is the same.

Each double coset $HgH$ determines:
\begin{itemize}
\item the \emph{intersection subgroup} $H \cap gHg^{-1} \leqslant H$
  (how much the two copies overlap),
\item the \emph{index} $[H : H \cap gHg^{-1}] = |HgH|/|H|$
  (how ``far apart'' the copies are).
\end{itemize}

\subsection{Twisted modules}

Let $M$ be a representation of $H$. For each $g \in G$, define $M_g$
to be the vector space $M$ with the $H$-action twisted by $g$: an
element $h \in H$ acts by
\[
  h \cdot_g m \;=\; (g^{-1}hg) \cdot m.
\]
This makes sense because $g^{-1}hg$ lies in $g^{-1}Hg \cong H$, so
we are just composing the original action with the automorphism
$h \mapsto g^{-1}hg$ of $H$. When $g \in H$, this is an inner
automorphism and $M_g \cong M$.

From $M_g$ we can build an $H$-module $N_g$ by restricting to the
intersection $H \cap gHg^{-1}$ and inducing back:
\[
  N_g \;:=\; \Ind_{H \cap gHg^{-1}}^{H}\;
  \Res_{H \cap gHg^{-1}}^{gHg^{-1}}\; M_g.
\]

\begin{lemma}\label{lem:Ng-well-defined}
If $g' = h_1 g h_2$ for some $h_1, h_2 \in H$, then
$N_{g'} \cong N_g$ as $H$-modules. So $N_g$ depends only on the
double coset $HgH$.
\end{lemma}

\begin{proof}
We check the two factors $h_1$ and $h_2$ separately.

\medskip\noindent\textbf{Right multiplication: $g' = gh_2$.}
Then $g'Hg'^{-1} = g(h_2 H h_2^{-1})g^{-1} = gHg^{-1}$ (since
$h_2 \in H$ normalises $H$). So the conjugate subgroup is unchanged,
the intersection $H \cap g'Hg'^{-1} = H \cap gHg^{-1}$ is unchanged,
and the twisted module $M_{g'} = M_{gh_2}$ has action
$h \cdot_{g'} m = (h_2^{-1} g^{-1} h g h_2) \cdot m
= h_2^{-1} \cdot (h \cdot_g (h_2 \cdot m))$.
This is $M_g$ conjugated by the action of $h_2 \in H$ on $M$, which
gives an isomorphic module after induction.

\medskip\noindent\textbf{Left multiplication: $g' = h_1 g$.}
Then $g'Hg'^{-1} = h_1(gHg^{-1})h_1^{-1}$, and
$H \cap g'Hg'^{-1} = H \cap h_1(gHg^{-1})h_1^{-1}
= h_1(H \cap gHg^{-1})h_1^{-1}$ (since $h_1 \in H$). Conjugation by
$h_1$ gives an isomorphism
\[
  \Ind_{H \cap gHg^{-1}}^{H}(\text{--})
  \;\xrightarrow{\;\sim\;}\;
  \Ind_{h_1(H \cap gHg^{-1})h_1^{-1}}^{H}(\text{--})
\]
(induction from conjugate subgroups gives isomorphic results), and
$M_{h_1 g}$ corresponds to $M_g$ under this conjugation. So
$N_{h_1 g} \cong N_g$.
\end{proof}

\subsection{The Mackey decomposition}

\begin{theorem}[Mackey]\label{thm:mackey}
Let $H \leqslant G$ and let $M$ be a representation of $H$. Then
\begin{equation}\label{eq:mackey-KH}
  \Res_H^G\;\Ind_H^G(M)
  \;\cong\;
  \bigoplus_{HgH \,\in\, H \backslash G / H} N_g
  \;=\;
  \bigoplus_{HgH \,\in\, H \backslash G / H}
  \Ind_{H \cap gHg^{-1}}^{H}\;
  \Res_{H \cap gHg^{-1}}^{gHg^{-1}}\;
  M_g.
\end{equation}
\end{theorem}

\begin{proof}
Recall that the induced representation is
\[
  \Ind_H^G(M) = k[G] \otimes_{k[H]} M,
\]
the tensor product of $M$ with the group algebra of $G$ over the group
algebra of $H$. The tensor is over $k[H]$, meaning we impose
$xh \otimes m = x \otimes h \cdot m$ for all $h \in H$. As a
$k$-vector space, choosing one representative $x$ from each left coset
$Hx \in H \backslash G$, we have
\[
  \Ind_H^G(M) = \bigoplus_{Hx \in H \backslash G} x \otimes M.
\]
An element $g \in G$ acts by left multiplication on the first factor:
\begin{equation}\label{eq:G-action-on-ind}
  g \cdot (x \otimes m) = (gx) \otimes m.
\end{equation}
If $gx$ lies in a different coset $Hy$, write $gx = hy$ for a unique
$h \in H$, and then $(gx) \otimes m = hy \otimes m = y \otimes h \cdot m$
(using the $k[H]$-tensor relation). So $g$ moves the element from the
$Hx$-summand to the $Hgx$-summand, and the $H$-part of the coset
change acts on $m$.

\begin{example}
$G$ acts on the set of left cosets $G/H = \{aH : a \in G\}$ by
left multiplication: $g \cdot (aH) = (ga)H$. This is well-defined
because $gaH = ga'H$ whenever $aH = a'H$. Linearising this action
(taking the vector space with basis $\{e_{aH}\}$ and extending
$g \cdot e_{aH} = e_{gaH}$ linearly) gives the \emph{permutation
representation} $k[G/H]$. The induced representation of the trivial
module is exactly this:
$\Ind_H^G(\mathbf{1}_H) \cong k[G/H]$.
\end{example}

Now consider $\Res_H^G \Ind_H^G(M)$: we have the $G$-module
$\Ind_H^G(M) = \bigoplus_{Hx \in H \backslash G} x \otimes M$,
and we restrict the action to $H \leqslant G$. The $G$-action
permuted the summands transitively, but $H$ only permutes them among
themselves according to how $H$ acts on $H \backslash G$. The coset
space $H \backslash G$ decomposes into orbits under the left
$H$-action, and these orbits are exactly the double cosets:
$H \backslash G = \bigsqcup_{HgH} (H \backslash HgH)$.
Within a single double coset $HgH$, the right cosets $\{Hx : x \in HgH\}$
form the set $H \backslash HgH$, which is in bijection with
$H / (H \cap gHg^{-1})$ via $Hx \leftrightarrow H \cap gHg^{-1} \cdot y$
(where $x = gy$ for appropriate $y$).

So the $H$-module $\Res_H^G \Ind_H^G(M)$ decomposes as
\[
  \Res_H^G \Ind_H^G(M)
  = \bigoplus_{HgH \in H \backslash G / H}
  \underbrace{\bigoplus_{Hx \in H \backslash HgH} x \otimes M}_{
    \text{contribution of the double coset } HgH}.
\]
It remains to identify each inner sum as
$\Ind_{H \cap gHg^{-1}}^{H}\, \Res_{H \cap gHg^{-1}}^{gHg^{-1}}\,
g \cdot M$.

Fix a double coset $HgH$. Any element $x \in HgH$ can be written as
$x = h g h'$ for some $h, h' \in H$. The right cosets of $H$ inside
$HgH$ are $\{Hg h' : h' \in H\}$, but $Hgh_1' = Hgh_2'$ if and only
if $h_2' (h_1')^{-1} \in g^{-1}Hg \cap H = g^{-1}(H \cap gHg^{-1})g$.
So the right cosets are indexed by
$(H \cap gHg^{-1}) \backslash H$ (from the right action of $H$ on $gH$).

Thus the inner sum, as an $H$-module under the left action, is
\[
  \bigoplus_{Hx \in H \backslash HgH} x \otimes M
  \;\cong\;
  k[H] \otimes_{k[H \cap gHg^{-1}]} (g \otimes M),
\]
where $H \cap gHg^{-1}$ acts on $g \otimes M$ by
$\gamma \cdot (g \otimes m) = g \otimes (g^{-1}\gamma g) \cdot m$
for $\gamma \in H \cap gHg^{-1}$ (noting that $g^{-1}\gamma g \in H$).
This is precisely
\[
  \Ind_{H \cap gHg^{-1}}^{H}\, \Res_{H \cap gHg^{-1}}^{gHg^{-1}}\,
  g \cdot M. \qedhere
\]
\end{proof}

\begin{remark}
More generally, for two subgroups $H, K \leqslant G$, the same
argument gives
\[
  \Res_K^G\;\Ind_H^G(M)
  \;\cong\;
  \bigoplus_{KgH \,\in\, K \backslash G / H}
  \Ind_{K \cap gHg^{-1}}^{K}\;
  \Res_{K \cap gHg^{-1}}^{gHg^{-1}}\;
  g \cdot M.
\]
The proof is identical, decomposing $K \backslash G$ into
$(K, H)$-double cosets instead of $(H, H)$-double cosets.
\end{remark}

Each double coset contributes one summand. The summand for $HgH$
routes $M$ through a \emph{bottleneck}: restrict to $H \cap gHg^{-1}$
(the overlap between the original and conjugate copies of $H$), then
induce back. The smaller the intersection, the tighter the bottleneck,
and the more the representation is ``reshuffled.''

\begin{itemize}
\item The \textbf{identity coset} $HeH = H$ gives intersection
  $H \cap H = H$ (no bottleneck). Its summand is $M$ itself.
\item Every \textbf{non-trivial coset} gives a proper subgroup
  $H \cap gHg^{-1} \lneq H$. The representation must squeeze through
  this smaller group and then expand back.
\end{itemize}

\begin{remark}\label{rem:G-disappears}
Notice that $G$ barely appears on the right-hand side
of~\eqref{eq:mackey-KH}. The sum is over the double cosets
$H \backslash G / H$ (which depend on $G$), and the conjugation
$g \cdot M$ uses an element $g \in G$. But every operation on
representations---the restriction, the induction, the subgroups
involved---takes place entirely within $H$ and its subgroups.
The ambient group $G$ sets the \emph{stage} (determining which double
cosets exist and which intersections arise), but does not
\emph{act} in the decomposition. The formula reduces a question about
$G$-representations to bookkeeping inside $H$.
\end{remark}

\subsection{The Hecke algebra}

\begin{definition}\label{def:hecke}
Let $\Hecke(G, H)$ be the space of $H$-bi-invariant functions on $G$:
\[
  \Hecke(G, H) = \{f : G \to k \mid f(h_1 g h_2) = f(g)
  \;\forall\, h_1, h_2 \in H\},
\]
with the \emph{convolution product}
$(f_1 * f_2)(g) = \sum_{x \in G} f_1(x)\, f_2(x^{-1}g)$.
\end{definition}

Since $H$-bi-invariant functions are constant on double cosets, the
characteristic functions $\{e_{HgH}\}$ form a basis, so
$\dim \Hecke(G, H) = |H \backslash G / H|$.

\begin{proposition}\label{prop:hecke-algebra}
$\Hecke(G, H)$ is a unital associative algebra under convolution.
\end{proposition}

\begin{proof}
We must check three things.

\medskip\noindent\textbf{Closure.} If $f_1, f_2$ are
$H$-bi-invariant, then so is $f_1 * f_2$. For any $h_1, h_2 \in H$:
\begin{align*}
  (f_1 * f_2)(h_1 g h_2)
  &= \sum_{x \in G} f_1(x)\, f_2(x^{-1} h_1 g h_2) \\
  &= \sum_{x \in G} f_1(x)\, f_2(x^{-1} h_1 g)
    \qquad\text{(right $H$-invariance of $f_2$)} \\
  &= \sum_{y \in G} f_1(y h_1)\, f_2(y^{-1} g)
    \qquad\text{(substituting $y = xh_1^{-1}$)} \\
  &= \sum_{y \in G} f_1(y)\, f_2(y^{-1} g)
    \qquad\text{(right $H$-invariance of $f_1$)} \\
  &= (f_1 * f_2)(g).
\end{align*}

\medskip\noindent\textbf{Associativity.} This follows from
associativity of multiplication in $G$:
\[
  ((f_1 * f_2) * f_3)(g)
  = \sum_{x,y} f_1(x)\, f_2(x^{-1}y)\, f_3(y^{-1}g)
  = (f_1 * (f_2 * f_3))(g).
\]

\medskip\noindent\textbf{Unit.} The function
$e_H = \frac{1}{|H|}\, \mathbf{1}_H$ (the characteristic function
of $H$, scaled by $1/|H|$) is the identity:
$(e_H * f)(g) = \frac{1}{|H|} \sum_{h \in H} f(h^{-1}g)
= f(g)$
by $H$-bi-invariance of $f$.
\end{proof}

We call $\Hecke(G, H)$ the \emph{Hecke algebra} of the pair $(G, H)$.

By the Frobenius--Mackey argument (Frobenius reciprocity applied
twice):
\begin{equation}\label{eq:hecke-end}
  \Hecke(G, H) \cong \End_G\bigl(\Ind_H^G(\mathbf{1}_H)\bigr).
\end{equation}
Each basis element $e_{HgH}$ is a $G$-endomorphism of the permutation
representation $k[G/H]$. In other words, every double coset defines an
operator on $G/H$, and these operators form an algebra under
composition.

%% ====================================================================
\section{Example 1: $S_n$ inside $S_{n+1}$}
\label{sec:Sn}
%% ====================================================================

\subsection{The embedding}

Embed $S_n$ into $S_{n+1}$ as the subgroup fixing the element $n{+}1$:
\[
  H = \Stab(n{+}1) = \{\sigma \in S_{n+1} : \sigma(n{+}1) = n{+}1\}
  \;\cong\; S_n.
\]
The coset space $S_{n+1}/S_n$ has $[S_{n+1}:S_n] = n+1$ elements,
one for each possible image of $n{+}1$: the coset $\sigma H$ consists
of all permutations that send $n{+}1$ to $\sigma(n{+}1)$.

\subsection{Conjugate copies of $S_n$}

A conjugate $gHg^{-1} = g\,\Stab(n{+}1)\,g^{-1} = \Stab(g(n{+}1))$
is the subgroup fixing $g(n{+}1)$. There are $n{+}1$ such subgroups:
\[
  \Stab(1),\; \Stab(2),\; \ldots,\; \Stab(n),\; \Stab(n{+}1) = H.
\]
All are copies of $S_n$ inside $S_{n+1}$, and they are all conjugate.

\subsection{Why there are exactly 2 double cosets}

From $H = \Stab(n{+}1)$'s perspective, the $n{+}1$ conjugate copies
fall into exactly \textbf{two} classes:

\begin{enumerate}
\item $\Stab(n{+}1)$ itself: the original copy. (1 subgroup.)

\item $\Stab(k)$ for $k \in \{1, \ldots, n\}$: a copy that fixes a
  \emph{different} point. These all look the same to $H$, because $H$
  acts transitively on $\{1, \ldots, n\}$---for any $k, k' \in
  \{1, \ldots, n\}$, there exists $h \in H$ with $h(k) = k'$, and
  conjugating $\Stab(k)$ by $h$ gives $\Stab(k')$. ($n$ subgroups,
  all equivalent under $H$-conjugacy.)
\end{enumerate}

This is the $H$-absorption principle: $n{+}1$ conjugate copies, but
$H$ can shuffle $n$ of them into each other, leaving 2 orbits.

The two double cosets are:
\begin{itemize}
\item \textbf{Identity coset} $HeH = H$: the set of permutations
  fixing $n{+}1$. Size $|H| = n!$.
\item \textbf{Non-trivial coset} $H(k,\,n{+}1)H$: the set of all
  permutations that \emph{move} $n{+}1$. Size
  $|S_{n+1}| - |H| = (n+1)! - n! = n \cdot n!$.
  (This is the same double coset for every $k \in \{1,\ldots,n\}$.)
\end{itemize}

\textbf{Dimension check}: $n! + n \cdot n! = n!(1 + n) = (n+1)!
= |S_{n+1}|$. \checkmark

\subsection{The intersection (the bottleneck)}

For the non-trivial double coset, pick representative $g = (k,\, n{+}1)$.
Then $gHg^{-1} = \Stab(k)$, and the intersection is:
\[
  H \cap gHg^{-1} = \Stab(n{+}1) \cap \Stab(k)
  = \{\sigma \in S_{n+1} : \sigma(n{+}1) = n{+}1 \text{ and } \sigma(k) = k\}
  \cong S_{n-1}.
\]
This is the subgroup that fixes \emph{both} points $n{+}1$ and $k$,
acting on the remaining $n-1$ elements.

\begin{remark}
The intersection $S_{n-1}$ depends on which representative $k$ we
choose, but different choices give conjugate copies of $S_{n-1}$ inside
$S_n$, and the Mackey summand is the same up to isomorphism.
\end{remark}

\subsection{The Mackey formula for $S_n \subset S_{n+1}$}

Combining the two double cosets, we get for any $S_n$-module $M$:
\begin{equation}\label{eq:mackey-Sn}
  \boxed{
  \Res_{S_n}^{S_{n+1}}\;\Ind_{S_n}^{S_{n+1}}(M)
  \;\cong\;
  M \;\oplus\; \Ind_{S_{n-1}}^{S_n}\,\Res_{S_{n-1}}^{S_n}(M).
  }
\end{equation}

\medskip\noindent\textbf{The identity coset} ($g = e$):
$H \cap eHe^{-1} = H$. The bottleneck is $H$ itself---nothing is
lost. The summand is $\Ind_H^H \Res_H^H(M) = M$. The representation
passes through unchanged.

\medskip\noindent\textbf{The non-trivial coset} ($g = (k,\, n{+}1)$):
$H \cap gHg^{-1} = S_{n-1}$. The bottleneck is one level smaller.
The summand routes $M$ through $S_{n-1}$:
\[
  M \;\xrightarrow{\;\Res\;}\;
  \Res_{S_{n-1}}^{S_n}(M) \;\xrightarrow{\;\Ind\;}\;
  \Ind_{S_{n-1}}^{S_n}\,\Res_{S_{n-1}}^{S_n}(M).
\]
The representation ``forgets'' one level of structure (restriction
to $S_{n-1}$) and then ``relearns'' it (induction back to $S_n$),
possibly ending up in a different irreducible component.

\begin{remark}[The Heisenberg relation]\label{rem:heisenberg}
Writing $U = \Ind_{S_n}^{S_{n+1}}$ and $D = \Res_{S_n}^{S_{n+1}}$,
equation~\eqref{eq:mackey-Sn} says
\[
  DU = \mathrm{Id} + UD,
\]
i.e., $DU - UD = \mathrm{Id}$. This is the \emph{Heisenberg
commutation relation}. The identity coset contributes the
$\mathrm{Id}$ and the non-trivial coset contributes $UD$.

Summing over all $n$ with generating function weights, this becomes
the Pieri operator identity
$\Omega^*[Xu]\,\Omega[Xv] = \frac{1}{1-uv}\,\Omega[Xv]\,\Omega^*[Xu]$,
where the factor $\frac{1}{1-uv}$ packages the identity coset
contributions across all levels.
\end{remark}

\subsection{What happens to an irreducible: walks on Young's lattice}

Let $M = S^\mu$ be the irreducible $S_n$-module labelled by a
partition $\mu \vdash n$. By the branching rule for the chain
$S_n \supset S_{n-1}$:

\medskip\noindent\textbf{Restriction} (remove one box):
$\Res_{S_{n-1}}^{S_n}(S^\mu) = \bigoplus_\nu S^\nu$,
summing over all partitions $\nu$ obtained from $\mu$ by removing
one box (i.e., $\mu/\nu$ is a single box).

\medskip\noindent\textbf{Induction} (add one box):
$\Ind_{S_{n-1}}^{S_n}(S^\nu) = \bigoplus_\lambda S^\lambda$,
summing over all partitions $\lambda \vdash n$ obtained from $\nu$
by adding one box.

\medskip\noindent Combining: the non-trivial Mackey summand is
\begin{equation}\label{eq:two-step-walk}
  \Ind_{S_{n-1}}^{S_n}\,\Res_{S_{n-1}}^{S_n}(S^\mu)
  = \bigoplus_\nu \bigoplus_{\lambda \supset \nu} S^\lambda,
\end{equation}
where $\nu$ ranges over $\mu$-minus-a-box and $\lambda$ ranges over
$\nu$-plus-a-box. This is a \textbf{two-step walk on Young's lattice}:
down one level and back up.

The complete Mackey decomposition is therefore:
\[
  \Res \circ \Ind(S^\mu) \;=\;
  S^\mu \;\oplus\;
  \bigoplus_{\substack{\nu \subset \mu \\ |\nu| = n-1}}
  \;\bigoplus_{\substack{\lambda \supset \nu \\ |\lambda| = n}}
  S^\lambda.
\]

\subsection{Worked example: $S_2$ inside $S_3$}

Take $n = 2$, so $H = S_2$ (fixing 3) inside $G = S_3$.
The two irreducibles of $S_2$ are $S^{(2)}$ (trivial) and
$S^{(1,1)}$ (sign).

\medskip\noindent\textbf{Case $M = S^{(2)}$}: Remove a box from
$(2)$: the only option is $\nu = (1)$. Add a box to $(1)$: get $(2)$
or $(1,1)$. So
\[
  \Res \circ \Ind(S^{(2)})
  = S^{(2)} \oplus \bigl(S^{(2)} \oplus S^{(1,1)}\bigr)
  = (S^{(2)})^{\oplus 2} \oplus S^{(1,1)}.
\]
Dimension check: $\dim = 2 + 1 = 3 = [S_3:S_2] \cdot \dim S^{(2)}
= 3 \cdot 1$. \checkmark

\medskip\noindent\textbf{Case $M = S^{(1,1)}$}: Remove a box from
$(1,1)$: the only option is $\nu = (1)$. Add a box to $(1)$: get
$(2)$ or $(1,1)$. So
\[
  \Res \circ \Ind(S^{(1,1)})
  = S^{(1,1)} \oplus \bigl(S^{(2)} \oplus S^{(1,1)}\bigr)
  = S^{(2)} \oplus (S^{(1,1)})^{\oplus 2}.
\]
Dimension check: $1 + 2 = 3$. \checkmark

\begin{remark}
In both cases, $\Res \circ \Ind(M)$ has dimension
$[G:H] \cdot \dim M = 3 \cdot 1 = 3$. The identity coset contributes
dimension 1 (returning $M$) and the non-trivial coset contributes
dimension 2 (the index $[S_2 : S_1] = 2$ times $\dim M = 1$).
\end{remark}

\subsection{The general case: $S_m \times S_n$ inside $S_{m+n}$}

Embed $H = S_m \times S_n$ into $G = S_{m+n}$ as the subgroup
preserving the partition $\{1, \ldots, m\} \mid \{m{+}1, \ldots, m{+}n\}$.
A conjugate $gHg^{-1}$ preserves a \emph{different} partition of
$\{1, \ldots, m{+}n\}$ into blocks of size $m$ and $n$.

From $H$'s perspective, what matters is: how much do the two
partitions overlap? If the new block of size $m$ shares exactly
$m - r$ elements with the old block of size $m$, then $r$ elements
``crossed over.''

\begin{proposition}
The double cosets $(S_m \times S_n) \backslash S_{m+n} / (S_m \times S_n)$
are indexed by $r = 0, 1, \ldots, \min(m,n)$, giving $\min(m,n) + 1$
double cosets. The intersection for the double coset indexed by $r$ is
\[
  H \cap gHg^{-1} \;\cong\;
  (S_{m-r} \times S_r) \times (S_r \times S_{n-r})
  \;\cong\;
  S_{m-r} \times S_r \times S_r \times S_{n-r},
\]
where $r$ elements from each block have crossed into the other.
\end{proposition}

The case $S_n \subset S_{n+1}$ is the special case $m = n$, $n = 1$
(or rather, $H = S_n \times S_1$, but $S_1$ is trivial). Then
$\min(n, 1) + 1 = 2$ double cosets, recovering our earlier result.

%% ====================================================================
\section{Example 2: $\GL_n(\F_q)$ and the Borel subgroup}
\label{sec:GLn}
%% ====================================================================

The right analogue of $S_n \subset S_{n+1}$ is not
$\GL_n(\F_q) \subset \GL_{n+1}(\F_q)$ (which has too many double
cosets for a clean story) but rather $(\GL_n(\F_q),\, B_n(\F_q))$,
where $B_n$ is the Borel subgroup of upper triangular matrices.

\subsection{The Borel subgroup}

Let $G = \GL_n(\F_q)$ and let $B = B_n(\F_q)$ be the subgroup of
upper triangular matrices:
\[
  B = \left\{\begin{pmatrix} a_{11} & a_{12} & \cdots & a_{1n} \\
  0 & a_{22} & \cdots & a_{2n} \\
  \vdots & & \ddots & \vdots \\
  0 & 0 & \cdots & a_{nn} \end{pmatrix} : a_{ii} \in \F_q^*,\;
  a_{ij} \in \F_q \text{ for } i < j\right\}.
\]
The order is $|B| = (q-1)^n \cdot q^{\binom{n}{2}}$.

\subsection{The flag variety}

$B$ is the stabiliser of the \emph{standard flag}
\[
  F_\bullet : \quad \{0\} \subset \langle e_1 \rangle \subset
  \langle e_1, e_2 \rangle \subset \cdots \subset
  \langle e_1, \ldots, e_n \rangle = \F_q^n.
\]
The coset space $G/B$ is the \emph{complete flag variety}: the set of
all complete flags $\{0\} = V_0 \subset V_1 \subset \cdots \subset
V_n = \F_q^n$ with $\dim V_i = i$.

For $n = 2$: a flag is just a line in $\F_q^2$, so $G/B =
\mathbb{P}^1(\F_q)$ with $q + 1$ points.

For $n = 3$: a flag is an incident line-plane pair $(\ell, \pi)$
with $\ell \subset \pi \subset \F_q^3$.

\subsection{The Bruhat decomposition}

The Weyl group of $\GL_n$ is $W = S_n$, realised as permutation
matrices. The key structural result:

\begin{theorem}[Bruhat]\label{thm:bruhat}
$\GL_n(\F_q) = \bigsqcup_{w \in S_n} BwB$. The double cosets
$B \backslash \GL_n / B$ are indexed by the symmetric group $S_n$.
\end{theorem}

\begin{proof}[Proof for $n = 2$]
$W = \{e, s\}$ where
$s = \bigl(\begin{smallmatrix} 0 & 1 \\ 1 & 0 \end{smallmatrix}\bigr)$.
Any $g = \bigl(\begin{smallmatrix} a & b \\ c & d \end{smallmatrix}\bigr)$:
if $c = 0$ then $g \in B = BeB$; if $c \neq 0$ then $g \in BsB$
(use row and column operations from $B$ to reduce $g$ to $s$). So
$\GL_2 = B \sqcup BsB$.
\end{proof}

Compare with the $S_n$ case: $S_n \backslash S_{n+1} / S_n$ had
$2$ double cosets (the point stabiliser absorbs everything). Here
$B \backslash \GL_n / B$ has $n!$ double cosets---one per permutation.
The Borel is ``too small'' to absorb conjugate copies, so the full
Weyl group is visible.

\subsection{The bottlenecks}

For each $w \in S_n$, the intersection $B \cap wBw^{-1}$ is the
bottleneck for the double coset $BwB$. Its size depends on $w$:
\begin{itemize}
\item $w = e$: $B \cap B = B$. Full overlap, no bottleneck.
\item $w = s_i$ (simple reflection): $B \cap s_i B s_i^{-1}$ has
  index $q$ in $B$---a thin bottleneck. The missing ``$q$ directions''
  correspond to the $q$ lines in a local $\mathbb{P}^1$.
\item $w = w_0$ (longest element): $B \cap w_0 B w_0^{-1} = T$
  (the torus, diagonal matrices). The tightest bottleneck---all the
  upper-triangular structure is lost.
\end{itemize}

In general, $[B : B \cap wBw^{-1}] = q^{\ell(w)}$ where $\ell(w)$ is
the Bruhat length (number of simple reflections in a reduced
expression for $w$). So longer elements have tighter bottlenecks.

\subsection{The Mackey formula for $(\GL_n, B)$}

Applying Theorem~\ref{thm:mackey} to the pair $(G, H) = (\GL_n, B)$:
\begin{equation}\label{eq:mackey-GLn-B}
  \Res_B^G\;\Ind_B^G(M)
  \;\cong\;
  \bigoplus_{w \in S_n}
  \Ind_{B \cap wBw^{-1}}^{B}\;
  \Res_{B \cap wBw^{-1}}^{wBw^{-1}}\; M_w.
\end{equation}
This has $n!$ summands, one per permutation. Each summand routes $M$
through the bottleneck $B \cap wBw^{-1}$, which has codimension
$\ell(w)$ in $B$.

\subsection{The quadratic relation from $\mathbb{P}^1(\F_q)$}

We now derive the Hecke algebra relations from this Mackey
decomposition. Start with $n = 2$, where $B \backslash \GL_2 / B$
has two elements and $\Hecke(\GL_2, B)$ is $2$-dimensional with
basis $\{T_e, T_s\}$.

The permutation representation $k[G/B] = k[\mathbb{P}^1(\F_q)]$
has $q + 1$ basis vectors $\delta_L$, one per line $L$ in $\F_q^2$.
The operator $T_s$ (the double coset $BsB$) sends each point to the
sum of all other points:
\begin{equation}\label{eq:Ts-action}
  T_s \cdot \delta_L = \sum_{L' \neq L} \delta_{L'}.
\end{equation}
This is because $BsB$ consists of matrices sending
$\langle e_1 \rangle$ to a different line, and $B$ acts transitively
on the $q$ lines other than $\langle e_1 \rangle$.

\begin{theorem}\label{thm:quadratic}
$(T_s - q)(T_s + 1) = 0$.
\end{theorem}

\begin{proof}
Apply $T_s$ twice to $\delta_L$:
\[
  T_s^2 \cdot \delta_L
  = \sum_{L' \neq L}\; \sum_{L'' \neq L'} \delta_{L''}.
\]
Count how often each $L''$ appears:
\begin{itemize}
\item $L'' = L$: each of the $q$ choices of $L' \neq L$ contributes
  $\delta_L$ once (since $L \neq L'$). Total: $q \cdot \delta_L$.
\item $L'' \neq L$: among the $q$ choices of $L'$ (all $\neq L$),
  exactly one is excluded ($L' = L''$), leaving $q - 1$.
  Total: $(q-1) \cdot \delta_{L''}$.
\end{itemize}
So $T_s^2 = q \cdot T_e + (q-1) \cdot T_s$, which factors as
$(T_s - q)(T_s + 1) = 0$.
\end{proof}

\begin{remark}
The eigenvalues $q$ and $-1$ correspond to:
$T_s = q$ on the trivial representation (the constant function:
$T_s$ sends it to $q$ times itself), and $T_s = -1$ on the
$q$-dimensional complement (functions summing to zero:
$T_s = J - I$ acts as $-I$ on $\ker J$).
\end{remark}

\subsection{The braid relation from flags in $\F_q^3$}

For $n = 3$, the generators $T_{s_1}$ and $T_{s_2}$ act on
flags $(\ell, \pi)$ by:
\begin{itemize}
\item $T_{s_1}$: replace the line $\ell$ by another line in $\pi$
  (sum over $q$ choices).
\item $T_{s_2}$: replace the plane $\pi$ by another plane through
  $\ell$ (sum over $q$ choices).
\end{itemize}
Each is a ``local $T_s$'' on a copy of $\mathbb{P}^1(\F_q)$
(the lines in a plane, or the planes through a line), so each
satisfies $(T_{s_i} - q)(T_{s_i} + 1) = 0$.

\begin{theorem}\label{thm:braid}
$T_{s_1} T_{s_2} T_{s_1} = T_{s_2} T_{s_1} T_{s_2}$.
\end{theorem}

\begin{proof}
Both sides produce a sum over flags at Bruhat distance $3$ from
$(\ell, \pi)$.
$T_{s_1}T_{s_2}T_{s_1}$ traces the path
$(\ell, \pi) \to (\ell', \pi) \to (\ell', \pi') \to (\ell'', \pi')$
(change line, change plane, change line).
$T_{s_2}T_{s_1}T_{s_2}$ traces
$(\ell, \pi) \to (\ell, \pi') \to (\ell', \pi') \to (\ell', \pi'')$
(change plane, change line, change plane).
Both equal $T_{w_0}$ (the basis element for the longest element of
$S_3$), because $s_1 s_2 s_1 = s_2 s_1 s_2 = w_0$ and both products
are length-additive: $\ell(w_0) = 3 = 1 + 1 + 1$.
\end{proof}

\subsection{The isomorphism}

\begin{definition}
The \emph{Iwahori--Hecke algebra} $H_n(q)$ is the associative algebra
generated by $T_1, \ldots, T_{n-1}$ subject to:
\begin{align}
  (T_i - q)(T_i + 1) &= 0
  && \text{for all } i, \label{eq:IH-quad} \\
  T_i T_{i+1} T_i &= T_{i+1} T_i T_{i+1}
  && \text{for all } i, \label{eq:IH-braid} \\
  T_i T_j &= T_j T_i
  && \text{for } |i - j| > 1. \label{eq:IH-comm}
\end{align}
At $q = 1$: $(T_i - 1)(T_i + 1) = 0$ gives $T_i^2 = 1$, and
$H_n(1) = \C[S_n]$.
\end{definition}

\begin{theorem}[Iwahori, 1964]\label{thm:iwahori}
$H_n(q) \cong \Hecke(\GL_n(\F_q),\, B_n(\F_q))$,
sending $T_i$ to the characteristic function of $B s_i B$.
Both algebras have dimension $n!$.
\end{theorem}

\subsection{Comparison: $S_n$ versus $(\GL_n, B)$}

\begin{center}
\renewcommand{\arraystretch}{1.5}
\begin{tabular}{@{}p{4cm}cc@{}}
\toprule
& $\boldsymbol{(S_{n+1},\, S_n)}$
& $\boldsymbol{(\GL_n(\F_q),\, B_n(\F_q))}$ \\
\midrule
Double cosets & $2$ & $n!$ (Bruhat) \\
Indexed by & \{identity, non-trivial\}
  & $S_n$ (Weyl group) \\
Hecke algebra dim & $2$ & $n!$ \\
Bottleneck for $w$ & $S_{n-1}$ & $B \cap wBw^{-1}$, index $q^{\ell(w)}$ \\
Quadratic relation & $T_s^2 = 1$ (involution)
  & $(T_s - q)(T_s + 1) = 0$ \\
At $q = 1$ & --- & recovers $\C[S_n]$ \\
Mackey gives & Heisenberg $DU - UD = I$
  & Iwahori--Hecke algebra $H_n(q)$ \\
Centralises & $\gl_N$ on $V^{\otimes n}$
  & $U_q(\gl_N)$ on $V^{\otimes n}$ \\
\bottomrule
\end{tabular}
\end{center}

The parameter $q$ counts the points of a ``local
$\mathbb{P}^1(\F_q)$'' (lines in a plane, or planes through a line).
The quadratic relation $(T_s - q)(T_s + 1) = 0$ encodes this
geometry: $T_s$ has eigenvalue $q$ on the constant function and $-1$
on its complement. As $q \to 1$, $\mathbb{P}^1$ shrinks to two
points, $T_s$ becomes a transposition, and $H_n(q)$ degenerates to
$\C[S_n]$.

%% ====================================================================
\section{The Hall algebra: $q$-deformed Heisenberg}
\label{sec:hall}
%% ====================================================================

For $S_n \subset S_{n+1}$, two double cosets gave the Heisenberg
relation $DU - UD = \mathrm{Id}$ (Remark~\ref{rem:heisenberg}). What
replaces this for $\GL_n(\F_q)$?

\subsection{Finite modules and their extensions}

Let $\F_q$ be a finite field. A \emph{finite $\F_q[t]$-module} is a
finite-dimensional $\F_q$-vector space $V$ equipped with an
$\F_q$-linear endomorphism $T : V \to V$. Equivalently, it is a
finitely generated $\F_q[t]$-module that is killed by some power of
$t$ (we consider only nilpotent endomorphisms, for simplicity).

By the structure theorem for modules over a PID, every such module is
isomorphic to
\[
  M_\lambda = \F_q[t]/(t^{\lambda_1}) \oplus \cdots \oplus
  \F_q[t]/(t^{\lambda_\ell}),
\]
for a unique partition $\lambda = (\lambda_1 \geqslant \cdots
\geqslant \lambda_\ell)$. The partition $\lambda$ records the sizes
of the Jordan blocks of $T$, and $|M_\lambda| = q^{|\lambda|}$.

\begin{definition}\label{def:hall-number}
For partitions $\lambda, \mu, \nu$, the \emph{Hall number}
$g^\lambda_{\mu\nu}(q)$ is the number of submodules
$N \subseteq M_\lambda$ such that $N \cong M_\nu$ and
$M_\lambda / N \cong M_\mu$:
\[
  g^\lambda_{\mu\nu}(q) = \bigl|\{N \leqslant M_\lambda :
  N \cong M_\nu,\; M_\lambda/N \cong M_\mu\}\bigr|.
\]
\end{definition}

The Hall numbers count \emph{extensions}: the number of ways to
build $M_\lambda$ as an extension of $M_\mu$ by $M_\nu$, up to the
automorphisms already accounted for. Philip Hall (1959) showed that
$g^\lambda_{\mu\nu}(q)$ is a polynomial in $q$ with integer
coefficients.

\begin{example}\label{ex:hall-number}
$g^{(2)}_{(1)(1)}(q)$: how many submodules $N \cong M_{(1)} = \F_q$
of $M_{(2)} = \F_q[t]/(t^2)$ have $M_{(2)}/N \cong M_{(1)}$? The
module $M_{(2)}$ has underlying space $\F_q^2$ with $T$ acting by
$T(a,b) = (b, 0)$. The only $T$-invariant subspace isomorphic to
$\F_q$ with the correct quotient is $\ker(T) = \{(a, 0)\}$.
So $g^{(2)}_{(1)(1)}(q) = 1$.

$g^{(1,1)}_{(1)(1)}(q)$: how many submodules $N \cong \F_q$ of
$M_{(1,1)} = \F_q \oplus \F_q$ (where $T = 0$) have quotient
$\cong \F_q$? Any one-dimensional subspace works, and there are
$q + 1$ of them (the lines in $\F_q^2$). But each is isomorphic to
$\F_q$ with $T = 0$, so $g^{(1,1)}_{(1)(1)}(q) = q + 1$.
\end{example}

\subsection{Definition of the Hall algebra}

\begin{definition}\label{def:hall-algebra}
The \emph{Hall algebra} $\mathcal{A}_q$ is the $\Z$-module freely
generated by symbols $[M_\lambda]$ for each partition $\lambda$, with
multiplication
\begin{equation}\label{eq:hall-mult}
  [M_\mu] \cdot [M_\nu] =
  \sum_\lambda g^\lambda_{\mu\nu}(q)\; [M_\lambda].
\end{equation}
\end{definition}

\begin{proposition}
The Hall algebra $\mathcal{A}_q$ is a unital associative algebra.
\end{proposition}

\begin{proof}
\emph{Associativity}: $([M_\mu] \cdot [M_\nu]) \cdot [M_\rho]$ counts
the number of chains $0 \subset N_1 \subset N_2 \subset M_\lambda$
with $N_1 \cong M_\rho$, $N_2/N_1 \cong M_\nu$,
$M_\lambda/N_2 \cong M_\mu$. This equals
$[M_\mu] \cdot ([M_\nu] \cdot [M_\rho])$, which counts the same chains
(a filtration of $M_\lambda$ with graded pieces $M_\rho, M_\nu, M_\mu$
is the same data regardless of whether we group the first two or last
two).

\emph{Unit}: $[M_\varnothing] = [0]$, corresponding to the empty
partition, satisfies $[0] \cdot [M_\lambda] = [M_\lambda]$ since the
only submodule isomorphic to $M_\lambda$ of $M_\lambda$ with zero
quotient is $M_\lambda$ itself.
\end{proof}

\subsection{The simplest case: multiplying two copies of $[M_{(1)}]$}

The element $[M_{(1)}]$ corresponds to the one-dimensional module
$\F_q$ (with $T = 0$). Computing $[M_{(1)}]^2$:
\[
  [M_{(1)}] \cdot [M_{(1)}]
  = g^{(2)}_{(1)(1)}(q)\, [M_{(2)}]
  + g^{(1,1)}_{(1)(1)}(q)\, [M_{(1,1)}].
\]
From Example~\ref{ex:hall-number}: $g^{(2)}_{(1)(1)} = 1$ and
$g^{(1,1)}_{(1)(1)} = q + 1$. So
\begin{equation}\label{eq:hall-square}
  [M_{(1)}]^2 = [M_{(2)}] + (q+1)\,[M_{(1,1)}].
\end{equation}
For comparison, in the ring of symmetric functions (the $q \to \infty$
/ $t \to 0$ limit), we have $h_1^2 = h_2 + e_2$, i.e., the
coefficient of $e_2$ is just $1$, not $q + 1$. The factor $q + 1$
counts the lines in $\F_q^2$---this is the finite field enrichment
that deforms the symmetric function product.

\subsection{The Hall algebra is a deformation of symmetric functions}

There is an isomorphism (the \emph{characteristic map})
\[
  \mathcal{A}_q \;\xrightarrow{\;\sim\;}\; \Lambda \otimes_\Z
  \Z[q],
\]
from the Hall algebra to the ring of symmetric functions with
coefficients extended to $\Z[q]$, sending $[M_\lambda]$ to a
scalar multiple of the Hall--Littlewood polynomial
$P_\lambda(x;\, q^{-1})$. Under this isomorphism, the Hall product
becomes a $q$-deformation of the product of symmetric functions.

In particular, defining $u_n = [M_{(1^n)}]$ (the class of
$\F_q^n$ with $T = 0$) and $e_n = [M_{(n)}]$ (the class of
$\F_q[t]/(t^n)$), we get $q$-deformed elementary and complete
symmetric functions.

\subsection{The commutation relation}

We now derive the $q$-Heisenberg relation. Define creation and
annihilation operators on $\mathcal{A}_q$ by:
\begin{align*}
  U(v) &= \text{multiplication by } \sum_{n \geqslant 0} [M_{(1^n)}]\, v^n
  \quad\text{(``add a semisimple module'')}, \\
  D(u) &= \text{adjoint operator (with respect to the Hall pairing)}.
\end{align*}
The Hall pairing is $\langle [M_\mu], [M_\nu] \rangle =
\delta_{\mu\nu}\, |\mathrm{Aut}(M_\mu)|$, where $|\mathrm{Aut}(M_\mu)|$
is the number of automorphisms of $M_\mu$.

The key computation: $D(u)\, U(v)$ counts pairs (submodule, quotient)
while $U(v)\, D(u)$ counts the same pairs in the opposite order. The
discrepancy between the two orderings is controlled by the number of
automorphisms, which introduces factors of $q$.

\begin{theorem}[Jing, 1991]\label{thm:HL-comm}
The vertex operators for the Hall algebra satisfy
\begin{equation}\label{eq:HL-comm}
  D(u)\, U(v) = \frac{1-uv/q}{1-uv}\; U(v)\, D(u).
\end{equation}
\end{theorem}

\begin{proof}[Proof sketch]
At the level of modes, the underlying Lie algebra relation is
\[
  [a_m, a_n^\dagger] = m\,(1 - q^{-|m|})\,\delta_{m+n,0},
\]
where $a_m, a_n^\dagger$ are the annihilation and creation modes.
At $q = \infty$ (i.e., $q^{-|m|} = 0$), this is the standard
Heisenberg relation $[a_m, a_n^\dagger] = m\,\delta_{m+n,0}$.

The factor $(1 - q^{-|m|})$ arises because the automorphism group
$\mathrm{Aut}(M_\lambda)$ has order divisible by powers of $q$: for
instance $|\mathrm{Aut}(\F_q^n)| = |\GL_n(\F_q)| =
\prod_{i=0}^{n-1}(q^n - q^i)$. When you commute creation past
annihilation, the ratio of automorphism groups for the two orderings
(submodule first vs.\ quotient first) produces the factor
$(1 - q^{-|m|})$.

Exponentiating the mode relation gives the generating function
form~\eqref{eq:HL-comm}.
\end{proof}

\subsection{Comparison with the Heisenberg relation}

Writing $t = 1/q$:

\begin{center}
\renewcommand{\arraystretch}{1.5}
\begin{tabular}{@{}p{5cm}p{6cm}@{}}
\toprule
\textbf{Symmetric groups ($t = 0$)} & \textbf{Hall algebra ($t = 1/q$)} \\
\midrule
$D(u)\, U(v) = \frac{1}{1-uv}\, U(v)\, D(u)$
  & $D(u)\, U(v) = \frac{1-tuv}{1-uv}\, U(v)\, D(u)$ \\
$[a_m, a_n^\dagger] = m\,\delta_{m+n,0}$
  & $[a_m, a_n^\dagger] = m(1 - t^{|m|})\,\delta_{m+n,0}$ \\
Schur functions $s_\lambda$
  & Hall--Littlewood $P_\lambda(x;\, t)$ \\
Hall numbers $= 0$ or $1$
  & Hall numbers $\in \Z[q]$ \\
$|\mathrm{Aut}(M_\lambda)| = 1$
  & $|\mathrm{Aut}(M_\lambda)| = $ polynomial in $q$ \\
\bottomrule
\end{tabular}
\end{center}

The extra factor $(1 - tuv)$ in the numerator can be traced to the
automorphism groups of $\F_q[t]$-modules. In the symmetric group
case ($q \to \infty$, $t \to 0$), these automorphism groups become
trivial and the factor disappears.

\subsection{Limiting cases}

\begin{itemize}
\item $t = 0$ ($q \to \infty$): the factor becomes $\frac{1}{1-uv}$,
  recovering the Heisenberg relation. The symmetric group case is the
  ``generic'' case where the finite field structure contributes
  nothing.

\item $t = 1$ ($q = 1$, formal limit): the factor becomes
  $\frac{1-uv}{1-uv} = 1$, so $DU = UD$ and creation and
  annihilation commute. This is the ``maximally abelian'' limit.

\item $t = 1/q$ for a prime power $q$: the actual $\GL_n(\F_q)$
  case. The commutation factor $\frac{1-uv/q}{1-uv}$ interpolates
  between the Heisenberg algebra and the abelian limit.
\end{itemize}

\begin{remark}
The Hall--Littlewood polynomials $P_\lambda(x;\, t)$ interpolate
between Schur functions ($t = 0$) and monomial symmetric functions
($t = 1$). At $t = 1/q$, they encode the representation theory of
$\GL_n(\F_q)$ in the same way that Schur functions encode $S_n$.
\end{remark}

%% ====================================================================
\section{The monad and its decomposition}
\label{sec:monad}
%% ====================================================================

\subsection{$\Ind \circ \Res$ as a monad}

For finite groups, induction and restriction participate in two
adjunctions:
\begin{align}
  \Ind_H^G &\;\dashv\; \Res_H^G
  && \text{(Frobenius reciprocity)},
  \label{eq:frob} \\
  \Res_H^G &\;\dashv\; \Ind_H^G
  && \text{(since $\Ind \cong \mathrm{Coind}$ for finite groups)}.
  \label{eq:res-left}
\end{align}
From adjunction~\eqref{eq:res-left}, $T = \Ind_H^G \circ \Res_H^G :
\Rep(G) \to \Rep(G)$ is a \textbf{monad}. The Mackey formula
decomposes this monad into a direct sum of functors, one per double
coset.

\subsection{The Hecke algebra and the monad}

The Hecke algebra $\Hecke(G, H)$ (Definition~\ref{def:hecke}) has
dimension $|H \backslash G / H|$: for $S_n \subset S_{n+1}$ this
is $2$, for $\GL_1(\F_3) \subset \GL_2(\F_3)$ this is $13$. By
the isomorphism $\Hecke(G,H) \cong \End_G(\Ind_H^G(\mathbf{1}_H))$
from~\eqref{eq:hecke-end}, each basis element $e_{HgH}$ is a
$G$-endomorphism of the permutation representation---and these are
exactly the adjacency operators of the Sabidussi coset graphs
$\Sab(G, H, g)$.

The monad perspective explains \emph{why} the Hecke algebra has a
multiplication: the monad $T = \Ind \circ \Res$ is composable, and
the monad multiplication $T^2 \to T$ (applying two Ind$\circ$Res
cycles and merging them) gives the convolution product on double
coset operators.

\subsection{Eilenberg--Moore algebras ($T$-algebras)}

A $T$-algebra is a $G$-representation $V$ with a structure map
$\alpha : T(V) = \Ind(\Res(V)) \to V$ satisfying unit and
associativity axioms. Via the Mackey decomposition, $\alpha$
decomposes as one projection $\alpha_{HgH}$ per double coset.

For finite groups in characteristic zero, the Barr--Beck theorem gives
$\Rep(G) \simeq T\text{-}\mathrm{Alg}$: a $G$-representation is the
same as an $H$-representation with coherent descent data (one
projection per double coset, satisfying compatibility conditions).

\subsection{The Kleisli category}

The Kleisli category $\Kl(T)$ has the same objects as $\Rep(G)$ but
morphisms
\[
  \Hom_{\Kl(T)}(V, W)
  = \Hom_G(V,\, \Ind\,\Res(W))
  \cong \Hom_H(\Res(V),\, \Res(W)).
\]
A Kleisli morphism is an $H$-equivariant map between restrictions:
$\Rep(G)$ with morphisms relaxed to $H$-equivariance.

The Hecke algebra appears as the Kleisli endomorphism algebra
of the trivial representation under the monad
$T' = \Res \circ \Ind$ on $\Rep(H)$:
\[
  \End_{\Kl(T')}(\mathbf{1}_H)
  \cong \Hom_G(\Ind(\mathbf{1}_H),\, \Ind(\mathbf{1}_H))
  = \Hecke(G, H).
\]
Each double coset = one Kleisli endomorphism = one Sabidussi graph
= one Mackey summand.

%% ====================================================================
\section{Connection to local rules and Schur--Weyl duality}
\label{sec:local-rules}
%% ====================================================================

\subsection{The Schur--Weyl double cosets}

For $G = S_{m+n}$ and $H = S_m \times S_n$, the $\min(m,n) + 1$
double cosets (indexed by the crossing parameter $r$) control the
Mackey decomposition of $\Ind \circ \Res$ for representations of
the symmetric group.

The Pieri operator $\Omega[Xv]$ (generating function for induction
along $S_n \hookrightarrow S_{n+1}$) and its dual $\Omega^*[Xu]$
satisfy
\[
  \Omega^*[Xu]\;\Omega[Xv]
  = \frac{1}{1-uv}\;\Omega[Xv]\;\Omega^*[Xu].
\]
This is the Mackey formula for the pair $(S_{n+1},\, S_n)$, with the
factor $\frac{1}{1-uv} = \sum_{m \geqslant 0}(uv)^m$ arising from
the non-trivial double coset.

\subsection{Local rules as Mackey steps}

Fomin's local rule for growth diagrams computes this Mackey
decomposition one cell at a time: the integer $m$ in the local rule
$\fU_{\alpha,\beta}(m, \mu) = \lambda$ is the identity-coset
contribution, and the partition data tracks the routing through the
non-trivial bottleneck.

\subsection{The categorical framework}

Davydov--Molev~\cite{DM} define Schur--Weyl categories
$\mathcal{C}(A_*)$ from multiplicative sequences of algebras.
The Hecke algebra appears in two guises:
\begin{enumerate}
\item $\Hecke(G, H)$: the double-coset Hecke algebra (this paper).
\item $\mathrm{H}_n(q)$: the Iwahori--Hecke algebra, which equals
  $\Hecke(\GL_n(\F_q),\, B_n(\F_q))$ for the Borel subgroup $B_n$.
\end{enumerate}
The passage from $k[S_n]$ to $\mathrm{H}_n(q)$ in Davydov--Molev
categorifies the passage from the symmetric group pair (2 double
cosets, clean bottleneck) to the $\GL_n(\F_q)$ pair (many double
cosets, field-arithmetic channels).

%% ====================================================================
\section{Summary}
\label{sec:summary}
%% ====================================================================

\begin{center}
\renewcommand{\arraystretch}{1.6}
\begin{tabular}{@{}p{3.2cm}p{4.2cm}p{4.8cm}@{}}
\toprule
\textbf{Concept}
  & $\boldsymbol{S_n \subset S_{n+1}}$
  & $\boldsymbol{\GL_1(\F_q) \subset \GL_2(\F_q)}$ \\
\midrule
Double cosets
  & 2 (fix same point, or different)
  & $q^2 + 2q - 2$ (field invariants) \\
Hecke algebra dim
  & 2
  & $q^2 + 2q - 2$ \\
Bottleneck
  & $S_{n-1}$ (one level down)
  & trivial (maximal reshuffling) \\
Central channels
  & none
  & $q - 2$ (from $Z(\GL_2)$) \\
Mackey formula
  & $M \oplus \Ind \Res(M)$
  & $(q{-}1)$ copies of $M$ plus reshuffled terms \\
Heisenberg relation
  & $UD = I + DU$
  & more complex commutation \\
Sabidussi graphs
  & 1 non-trivial graph (valency $n$)
  & $q^2 + 2q - 3$ non-trivial graphs \\
\bottomrule
\end{tabular}
\end{center}

The unifying principle: double cosets count how many essentially
different ways a conjugate copy of $H$ can sit inside $G$, up to the
symmetry of $H$ itself. For $S_n$, $H$'s transitivity absorbs almost
all the copies, leaving just 2 classes. For $\GL_n(\F_q)$, the field
structure creates additional invariants ($d$, the normalised $A$-entry)
that $H$ cannot absorb, producing polynomially many classes in $q$.
The Mackey formula decomposes the monad $\Ind \circ \Res$ along these
classes, the Hecke algebra records their convolution structure, and the
local rules / Sabidussi graphs are individual channels of the
decomposition.

\begin{thebibliography}{99}

\bibitem{DM}
A.~Davydov and A.~Molev,
\emph{A categorical approach to classical and quantum Schur--Weyl
duality},
preprint.

\bibitem{Fomin}
S.~Fomin,
\emph{Schensted algorithms for dual graded graphs},
J.~Algebraic Combin.\ \textbf{4} (1995), 5--45.

\bibitem{Iwahori}
N.~Iwahori,
\emph{On the structure of a Hecke ring of a Chevalley group over a
finite field},
J.~Fac.\ Sci.\ Univ.\ Tokyo Sect.\ I \textbf{10} (1964), 215--236.

\bibitem{Langer}
R.~Langer,
\emph{Cylindric partitions and Littlewood--Richardson coefficients},
PhD thesis, University of Melbourne, 2012.

\bibitem{Mackey}
G.~W.~Mackey,
\emph{On induced representations of groups},
Amer.\ J.\ Math.\ \textbf{73} (1951), 576--592.

\bibitem{Sabidussi}
G.~Sabidussi,
\emph{Vertex-transitive graphs},
Monatsh.\ Math.\ \textbf{68} (1964), 426--438.

\end{thebibliography}

\end{document}
